\section{引言}
\subsection{问题背景}

随着人工智能与云计算规模持续扩张，数据中心已成为电力系统中不可忽视的负荷主体，其高密度、强弹性、跨区域的特征使得"算力"与"电力"的耦合日益加深。国家"东数西算"工程将大量计算任务部署至西部富可再生能源地区，为清洁算力提供了新路径；而另一方面，风电、光伏出力具有强波动性与不可调度性，若算力负荷不能随新能源出力联动调整，将产生大量弃风、弃光，或被迫高价购电、高碳排运行。因此，如何在满足任务服务质量（时延、截止、GPU 容量）的前提下，把算力调度与电力调度联合优化（算电协同），从而同时降低运行成本与碳排放，是兼具理论与工程价值的研究课题。

本文所研究的问题取自某算力网络的真实运行样例：包含六个区域（A/B/C 为东部算力中心，D 为西部算力中心，E 为光伏富集区，F 为风电富集区），50000 条非抢占任务运行在 0--2399 小时到达、2406 小时前全部完成，并给出逐时电价、碳强度、可用新能源、储能参数与任务 GPU—功率映射等配套数据。核心矛盾在于：任务调度改变各区域逐时 IT 负荷，进而改变购电、外送、弃电与碳排放；而储能又能进一步平抑负荷与新能源的失配。四个子问题由此递进展开。

\subsection{问题重述}

\textbf{问题1：} 对任务与电力数据进行统计分析和未来 24 小时到达 GPU 需求的预测；在不考虑储能的前提下，给出满足时延与 GPU 容量约束的基础调度方案，并分析目标值、GPU 使用效率，给出最后 24 小时（2376--2406）的收尾调度计划与甘特图。

\textbf{问题2：} 在引入碳约束、考虑购电、外送与弃电机制、不使用储能的条件下，从任务柔性与电网交互两个角度建立碳—成本最优的任务调度模型，得到公平可比的调度方案，并给出碳排放约束下的成本—效益分析。

\textbf{问题3：} 在不需要跨区域调度、使用储能、使用部分或全部附件负荷（含非 AI 负荷）固定不变的条件下，建立储能协同优化模型，求解购售电与充放电策略。

\textbf{问题4：} 对系统进行碳—成本联合优化，同时优化任务时序与储能充放电，构建联合调度模型；设置情景（新能源出力降低、电价机制变化等）检验鲁棒性，并估计"仅任务侧""仅储能侧""双侧联合"三者的最优边界与交互关系。
